% ============================================================
% 测试卷 — 第4章 指数函数与对数函数
% ============================================================
\documentclass[12pt, a4paper]{ctexart}
\usepackage[margin=2cm]{geometry}
\usepackage{amsmath, amssymb}
\usepackage{enumitem}
\usepackage{fancyhdr}
\usepackage{xcolor}
\usepackage{tikz}
\usepackage{graphicx}
\usepackage{tabularx}
\usepackage{multicol}

\definecolor{mainblue}{RGB}{41, 98, 168}

\pagestyle{fancy}
\fancyhf{}
\fancyhead[L]{\small\textcolor{mainblue}{2025-2026学年度·第4章测试卷}}
\fancyhead[R]{\small\textcolor{mainblue}{数学（必修第一册）}}
\fancyfoot[C]{\small —\ \thepage\ —}
\renewcommand{\headrulewidth}{0.6pt}
\renewcommand{\headrule}{\hbox to\headwidth{\color{mainblue}\leaders\hrule height \headrulewidth\hfill}}

\setCJKmainfont{Hiragino Sans GB}
\setCJKsansfont{Heiti SC}

\newcommand{\blank}[1]{\underline{\hspace{#1}}}

\begin{document}

\begin{center}
{\Large\bfseries 第4章\quad 指数函数与对数函数\quad 测试卷}\\[6pt]
{\normalsize 数学（必修第一册）}\\[4pt]
{\small 考试时间：120分钟\quad 满分：150分}
\end{center}

\vspace{2pt}
\begin{center}
\begin{tabular}{|c|c|c|c|c|}
\hline
题号 & 一 & 二 & 三 & 总分 \\ \hline
得分 & & & & \\ \hline
\end{tabular}
\end{center}
\vspace{4pt}

\noindent\textbf{注意事项：}
\begin{enumerate}[leftmargin=2em]
  \item 答题前，请在答题卡上填写姓名、班级和考号。
  \item 选择题请用2B铅笔在答题卡上涂黑对应选项。
  \item 非选择题请用黑色签字笔在答题卡指定区域作答。
\end{enumerate}

\vspace{6pt}

% ============ 一、单项选择题 ============
\noindent\rule{\textwidth}{0.4pt}\\[2pt]
{\large\bfseries 一、单项选择题}（本大题共8小题，每小题5分，共40分。每小题只有一个选项符合题目要求。）

\vspace{6pt}

\begin{enumerate}
\item 计算 $2^{\frac{1}{2}}\cdot 2^{\frac{3}{2}}$ 的值为（\quad）
  \begin{enumerate}[label=\Alph*., itemsep=0pt]
    \item $2$ \quad \item $4$ \quad \item $2\sqrt{2}$ \quad \item $8$
  \end{enumerate}

\item 函数 $y=3^x$ 的图象必过点（\quad）
  \begin{enumerate}[label=\Alph*., itemsep=0pt]
    \item $(0,0)$ \quad \item $(0,1)$ \quad \item $(1,0)$ \quad \item $(1,1)$
  \end{enumerate}

\item $\log_2 8$ 的值为（\quad）
  \begin{enumerate}[label=\Alph*., itemsep=0pt]
    \item $2$ \quad \item $3$ \quad \item $4$ \quad \item $8$
  \end{enumerate}

\item 函数 $f(x)=\log_2(x-1)$ 的定义域为（\quad）
  \begin{enumerate}[label=\Alph*., itemsep=0pt]
    \item $(0,+\infty)$ \quad \item $(1,+\infty)$ \quad \item $[1,+\infty)$ \quad \item $(2,+\infty)$
  \end{enumerate}

\item 下列不等式正确的是（\quad）
  \begin{enumerate}[label=\Alph*., itemsep=0pt]
    \item $2^{0.3}<2^{0.2}$ \quad \item $0.3^2>0.3^{0.5}$
    \item $\log_2 3>\log_2 5$ \quad \item $\log_{0.5} 3<\log_{0.5} 2$
  \end{enumerate}

\item 函数 $y=\left(\dfrac{1}{2}\right)^x$ 在 $\mathbb{R}$ 上是（\quad）
  \begin{enumerate}[label=\Alph*., itemsep=0pt]
    \item 单调递增函数 \quad \item 单调递减函数
    \item 先增后减的函数 \quad \item 先减后增的函数
  \end{enumerate}

\item 若 $\log_a 2<\log_a 3$，则 $a$ 的取值范围为（\quad）
  \begin{enumerate}[label=\Alph*., itemsep=0pt]
    \item $a>1$ \quad \item $0<a<1$
    \item $a>0$ \quad \item $a>2$
  \end{enumerate}

\item 函数 $f(x)=a^{x-1}+1$（$a>0$ 且 $a\neq 1$）的图象恒过定点（\quad）
  \begin{enumerate}[label=\Alph*., itemsep=0pt]
    \item $(0,2)$ \quad \item $(1,1)$ \quad \item $(1,2)$ \quad \item $(0,1)$
  \end{enumerate}
\end{enumerate}

\vspace{6pt}

% ============ 二、多项选择题 ============
\noindent\rule{\textwidth}{0.4pt}\\[2pt]
{\large\bfseries 二、多项选择题}（本大题共4小题，每小题5分，共20分。全部选对得5分，部分选对得2分，有选错的得0分。）

\vspace{6pt}

\begin{enumerate}[resume]
\item 下列运算正确的是（\quad）
  \begin{enumerate}[label=\Alph*., itemsep=0pt]
    \item $a^m\cdot a^n=a^{mn}$
    \item $(a^m)^n=a^{mn}$
    \item $(ab)^n=a^n b^n$
    \item $\dfrac{a^m}{a^n}=a^{m-n}$（$a\neq 0$）
  \end{enumerate}

\item 关于函数 $y=\ln x$，下列说法正确的是（\quad）
  \begin{enumerate}[label=\Alph*., itemsep=0pt]
    \item 定义域为 $(0,+\infty)$
    \item 值域为 $\mathbb{R}$
    \item 在定义域上单调递增
    \item 过点 $(1,0)$
  \end{enumerate}

\item 已知 $a=\log_3 2$，$b=\log_5 3$，$c=2^{0.5}$，则（\quad）
  \begin{enumerate}[label=\Alph*., itemsep=0pt]
    \item $0<a<1$
    \item $0<b<1$
    \item $c>1$
    \item $a<b$
  \end{enumerate}

\item 关于函数 $f(x)=2^x-2^{-x}$，下列结论正确的是（\quad）
  \begin{enumerate}[label=\Alph*., itemsep=0pt]
    \item $f(x)$ 是奇函数
    \item $f(x)$ 在 $\mathbb{R}$ 上单调递增
    \item $f(x)$ 的值域为 $\mathbb{R}$
    \item $f(1)>1$
  \end{enumerate}
\end{enumerate}

\vspace{6pt}

% ============ 三、填空题 ============
\noindent\rule{\textwidth}{0.4pt}\\[2pt]
{\large\bfseries 三、填空题}（本大题共4小题，每小题5分，共20分。）

\vspace{6pt}

\begin{enumerate}[resume]
\item 化简 $8^{\frac{2}{3}}$ 的值为 \underline{\hspace{2cm}}。

\item 方程 $2^x=16$ 的解为 $x=$ \underline{\hspace{2cm}}。

\item 已知 $\lg 2\approx 0.3010$，$\lg 3\approx 0.4771$，则 $\lg 12\approx$ \underline{\hspace{3cm}}。

\item 函数 $f(x)=\log_2(x^2-3x+2)$ 的定义域为 \underline{\hspace{4cm}}。
\end{enumerate}

\vspace{6pt}

% ============ 四、解答题 ============
\noindent\rule{\textwidth}{0.4pt}\\[2pt]
{\large\bfseries 四、解答题}（本大题共6小题，共70分。解答应写出文字说明、证明过程或演算步骤。）

\vspace{6pt}

\begin{enumerate}[resume]
\item[17.]（本小题10分）\\
计算下列各式的值：
\begin{enumerate}[(1)]
  \item $\left(\dfrac{8}{27}\right)^{-\frac{2}{3}}+\log_3 9+\lg 100$；
  \item $2\log_3 2-\log_3 \dfrac{32}{9}+\log_3 8$。
\end{enumerate}

\vspace{6cm}

\item[18.]（本小题12分）\\
已知函数 $f(x)=a^x$（$a>0$ 且 $a\neq 1$）的图象经过点 $(2,9)$。
\begin{enumerate}[(1)]
  \item 求 $a$ 的值；
  \item 求 $f(-1)$、$f(0)$、$f\!\left(\dfrac{1}{2}\right)$ 的值；
  \item 比较 $f(-2)$ 与 $f(-3)$ 的大小。
\end{enumerate}

\vspace{6cm}

\item[19.]（本小题12分）\\
求下列方程的解：
\begin{enumerate}[(1)]
  \item $4^x-3\cdot 2^x+2=0$；
  \item $\log_2(x+1)+\log_2(x-1)=3$。
\end{enumerate}

\vspace{6cm}

\item[20.]（本小题12分）\\
已知函数 $f(x)=\log_a(x+1)$（$a>0$ 且 $a\neq 1$）。
\begin{enumerate}[(1)]
  \item 求 $f(x)$ 的定义域；
  \item 当 $a=2$ 时，求使 $f(x)>0$ 成立的 $x$ 的取值范围；
  \item 当 $a>1$ 时，证明 $f(x)$ 在定义域上单调递增。
\end{enumerate}

\vspace{6cm}

\item[21.]（本小题12分）\\
已知函数 $f(x)=2^x-2^{-x}$。
\begin{enumerate}[(1)]
  \item 判断 $f(x)$ 的奇偶性，并说明理由；
  \item 判断 $f(x)$ 在 $\mathbb{R}$ 上的单调性，并用定义证明。
\end{enumerate}

\vspace{7cm}

\item[22.]（本小题12分）\\
某城市人口的年增长率为 $2\%$，2020年底该城市人口为 $100$ 万。设经过 $x$ 年后（$x\geq 0$，$x\in\mathbb{N}$），该城市人口为 $y$ 万。
\begin{enumerate}[(1)]
  \item 写出 $y$ 关于 $x$ 的函数关系式；
  \item 按此增长率，大约经过多少年后该城市人口超过 $120$ 万？（参考数据：$\lg 1.02\approx 0.0086$，$\lg 1.2\approx 0.0792$）
\end{enumerate}

\vspace{7cm}
\end{enumerate}

\vfill
\begin{center}
\small\textcolor{gray}{本试卷仅供教学参考}
\end{center}

\end{document}
